\documentclass[11pt]{article}

\usepackage[margin=1in]{geometry}
\usepackage{booktabs}

\title{Lab 1: Image Sharpening \\ \large Operating Systems}
\author{Ayush \\ CS24BT001 \\ IIT Dharwad}
\date{\today}

\begin{document}
\maketitle

\section*{Objective}
Implement an image sharpening pipeline in C++ using unsharp masking, reading and
writing PPM images with the given \texttt{libppm} functions, and measure how long
the pipeline takes to run.

\section*{Approach}
Unsharp masking works in three steps:

\begin{enumerate}
	\item \textbf{Smoothen} -- replace every pixel with the average of its
	3$\times$3 neighborhood. Pixels on the border just average over whatever
	neighbors they actually have.
	\item \textbf{Find details} -- subtract the smoothened image from the original
	to get the high-frequency detail. Since pixel values are unsigned, negative
	differences are clamped to 0 instead of subtracting directly.
	\item \textbf{Sharpen} -- add the detail back onto the original image, clamped
	at 255, to get the final sharpened output.
\end{enumerate}

While writing this, I initially summed the neighborhood directly into a
\texttt{uint8\_t}, which overflows since 9 pixels can sum to more than 255 -- fixed
by accumulating in a wider integer type before dividing. I also ran into the
unsigned subtraction issue mentioned above, since a plain \texttt{input -
smoothened} wraps around whenever a pixel is darker than its neighborhood average,
which happens all over the image and not just at edges.


Timing is measured around \texttt{read\_ppm\_file}, S1, S2 and S3 using
\texttt{std::chrono::high\_resolution\_clock}, and logged (along with the image
number) both to the console and to \texttt{timing.txt}.

\section*{Timing Results}
Ran the program on 7 test images of increasing size using
\texttt{make run-sharpen INPUT=n OUTPUT=no}. Results from \texttt{timing.txt}:

\begin{table}[h]
	\centering
	\begin{tabular}{cr}
		\toprule
		Image & Time (microseconds) \\
		\midrule
		1 & 25{,}357 \\
		2 & 105{,}165 \\
		3 & 165{,}486 \\
		4 & 310{,}888 \\
		5 & 410{,}679 \\
		6 & 891{,}836 \\
		7 & 2{,}426{,}501 \\
		\bottomrule
	\end{tabular}
\end{table}

\section*{Analysis}
The bigger the input image, the longer the program takes to run -- makes sense
since every extra pixel adds a fixed amount of work across the three stages.

\section*{Conclusion}
The pipeline smooths, extracts detail, and sharpens PPM images correctly once the
overflow and underflow bugs in the pixel arithmetic were fixed.

\end{document}